\chapter{Justificaci\'on detallada de decisiones}\label{justificacion}

\section{Introducci\'on}

En un TFG de ingenier\'ia de software, implementar funcionalidades no es suficiente
si no se explica por qu\'e se ha elegido cada opci\'on t\'ecnica,
qu\'e alternativas se consideraron y qu\'e compromisos se asumieron.

Este cap\'itulo ampl\'ia de forma exhaustiva la justificaci\'on de decisiones
adoptadas durante el proyecto,
con foco en cuatro ejes: producto, arquitectura, calidad y operaci\'on.

\section{Marco de toma de decisiones}

\subsection{Criterios utilizados}

La selecci\'on de alternativas se realiz\'o usando un enfoque multicriterio,
priorizando los siguientes factores:

\begin{itemize}
  \item alineaci\'on con requisitos funcionales y no funcionales del ERS;
  \item viabilidad temporal dentro de 660 horas de proyecto;
  \item mantenibilidad y capacidad de evoluci\'on posterior;
  \item coste operativo compatible con entorno acad\'emico;
  \item seguridad y trazabilidad del ciclo completo.
\end{itemize}

\subsection{Ponderaci\'on de criterios}

Para evitar decisiones ad-hoc,
se aplic\'o una ponderaci\'on orientativa de criterios
que se mantuvo estable en los principales bloques de dise\~no.

\begin{table}[ht]
\centering
\begin{tabular}{|P{6.4cm}|P{2.5cm}|}
\hline
\textbf{Criterio} & \textbf{Peso (\%)} \\
\hline
Adecuaci\'on funcional y cobertura de requisitos & 30 \\
\hline
Mantenibilidad y evoluci\'on & 20 \\
\hline
Tiempo de implementaci\'on y riesgo de entrega & 20 \\
\hline
Coste operativo y sostenibilidad & 15 \\
\hline
Seguridad y robustez t\'ecnica & 15 \\
\hline
\textbf{Total} & \textbf{100} \\
\hline
\end{tabular}
\caption{Ponderaci\'on orientativa para selecci\'on de decisiones}
\label{tab:pesos-criterios}
\end{table}

\section{Decisiones de alcance funcional}

\subsection{Decisi\'on D1: definir un MVP centrado en flujo acad\'emico nuclear}

\textbf{Alternativas consideradas}:
\begin{itemize}
  \item incorporar desde inicio todas las funcionalidades del backlog;
  \item limitar el alcance a un MVP con ciclo completo de entrega y evaluaci\'on.
\end{itemize}

\textbf{Decisi\'on adoptada}: segunda alternativa.

\textbf{Raz\'on principal}:
maximizar valor entregado en plazo,
evitando dispersi\'on en funcionalidades accesorias
que pod\'ian comprometer la estabilidad del n\'ucleo.

\textbf{Impacto positivo}:
\begin{itemize}
  \item mayor probabilidad de cierre funcional completo;
  \item mejor cobertura de pruebas sobre procesos cr\'iticos;
  \item documentaci\'on m\'as s\'olida de extremo a extremo.
\end{itemize}

\textbf{Compromiso asumido}:
funcionalidades \emph{Should/Could} pasan a hoja de ruta,
sin bloquear la validez del resultado principal.

\subsection{Decisi\'on D2: dise\~nar experiencia por rol desde el inicio}

\textbf{Decisi\'on adoptada}:
separar claramente experiencias de administrador, profesor y estudiante.

\textbf{Razones}:
\begin{itemize}
  \item reduce ambig\"uedad de permisos en interfaz y API;
  \item simplifica pruebas de autorizaci\'on y casos negativos;
  \item mejora usabilidad al mostrar solo operaciones pertinentes.
\end{itemize}

\section{Decisiones de arquitectura}

\subsection{Decisi\'on D3: arquitectura SPA + API desacoplada}

\textbf{Alternativas}:
\begin{itemize}
  \item aplicaci\'on monol\'itica renderizada en servidor;
  \item frontend SPA + backend REST desacoplado.
\end{itemize}

\textbf{Decisi\'on adoptada}: SPA + API.

\textbf{Razones de peso}:
\begin{itemize}
  \item separaci\'on de responsabilidades y menor acoplamiento;
  \item evoluci\'on independiente de capa de presentaci\'on y l\'ogica;
  \item mejor ajuste a un dominio con flujos ricos de interfaz (evaluaciones, filtros, historial).
\end{itemize}

\textbf{Trade-off}:
aumenta esfuerzo de integraci\'on y de pruebas E2E,
compensado por mayor mantenibilidad a medio plazo.

\subsection{Decisi\'on D4: arquitectura backend por capas}

\textbf{Decisi\'on adoptada}:
patr\'on \texttt{controller-service-repository-entity} con DTOs de frontera.

\textbf{Razones}:
\begin{itemize}
  \item control claro de responsabilidades;
  \item testabilidad por m\'odulo;
  \item evita exponer modelo de persistencia en contratos externos;
  \item facilita versionado de API y evoluci\'on de dominio.
\end{itemize}

\subsection{Decisi\'on D5: uso de DTOs y mapper dedicado}

\textbf{Problema que resuelve}:
evitar dependencia directa entre entidades JPA y respuestas HTTP.

\textbf{Razones}:
\begin{itemize}
  \item minimiza acoplamiento entre capa de datos y capa API;
  \item reduce riesgo de sobreexposici\'on de atributos sensibles;
  \item permite modelar respuestas orientadas al caso de uso.
\end{itemize}

\textbf{Coste asumido}:
mayor volumen de c\'odigo estructural,
aceptado por mejora en mantenibilidad y seguridad de contrato.

\section{Decisiones de stack tecnol\'ogico}

\subsection{Decisi\'on D6: backend en Spring Boot + Java 21}

\textbf{Razones principales}:
\begin{itemize}
  \item ecosistema maduro en seguridad, validaci\'on y persistencia;
  \item integraci\'on robusta con Spring Security, JPA y testing;
  \item soporte industrial y documentaci\'on extensa.
\end{itemize}

\textbf{Comparativa resumida}:
frente a opciones como Node.js o Django,
Spring reduce incertidumbre de ensamblaje en aspectos cr\'iticos
(seguridad por rol, transacciones, control de acceso por contexto).

\subsection{Decisi\'on D7: frontend en React + TypeScript + Vite}

\textbf{Razones}:
\begin{itemize}
  \item equilibrio entre productividad y control tipado;
  \item ecosistema amplio para rutas, formularios y testing;
  \item arranque y build r\'apidos con Vite.
\end{itemize}

\textbf{Compromiso}:
mayor libertad implica mayor disciplina de arquitectura en frontend,
mitigada mediante organizaci\'on por m\'odulos (pages/components/context/services).

\section{Decisiones de seguridad}

\subsection{Decisi\'on D8: autenticaci\'on JWT con refresh token}

\textbf{Motivaci\'on}:
compatible con arquitectura desacoplada y despliegue stateless.

\textbf{Razones de elecci\'on}:
\begin{itemize}
  \item reduce estado de sesi\'on en servidor;
  \item habilita renovaci\'on transparente de sesi\'on en cliente;
  \item facilita escalado horizontal del backend.
\end{itemize}

\textbf{Riesgos identificados}:
control de expiraci\'on y revocaci\'on m\'as complejo que sesi\'on tradicional.

\textbf{Mitigaci\'on aplicada}:
expiraciones diferenciadas, refresh controlado,
filtros de seguridad y validaci\'on de contexto por recurso.

\subsection{Decisi\'on D9: autorizaci\'on por rol + contexto acad\'emico}

\textbf{Raz\'on}:
un rol global no es suficiente en un dominio donde un usuario
puede operar con permisos distintos seg\'un curso, grupo o actividad.

\textbf{Aplicaci\'on concreta}:
validaci\'on en controladores y servicios de pertenencia real
al contexto objetivo antes de permitir operaciones sensibles.

\section{Decisiones de datos e integraciones}

\subsection{Decisi\'on D10: PostgreSQL como base de datos principal}

\textbf{Razones}:
\begin{itemize}
  \item robustez transaccional y consistencia;
  \item buena compatibilidad con JPA/Hibernate;
  \item adaptaci\'on sencilla a entornos gestionados.
\end{itemize}

\subsection{Decisi\'on D11: Flyway como autoridad de esquema}

\textbf{Razones}:
\begin{itemize}
  \item trazabilidad expl\'icita de cambios en base de datos;
  \item despliegues reproducibles entre entornos;
  \item menor riesgo de divergencias silenciosas de esquema.
\end{itemize}

\textbf{Decisi\'on asociada}:
usar \texttt{ddl-auto=validate} en producci\'on para evitar
mutaciones impl\'icitas del esquema fuera de control.

\subsection{Decisi\'on D12: estrategia h\'ibrida de almacenamiento de archivos}

\textbf{Decisi\'on adoptada}:
almacenamiento local en desarrollo y cloud persistente en producci\'on.

\textbf{Razones}:
\begin{itemize}
  \item simplicidad y velocidad en entorno local;
  \item persistencia y resiliencia en entornos con filesystem ef\'imero;
  \item capacidad de habilitar integraciones externas sin romper n\'ucleo.
\end{itemize}

\section{Decisiones de calidad y verificaci\'on}

\subsection{Decisi\'on D13: testing en capas}

\textbf{Razones}:
\begin{itemize}
  \item unitarias para retroalimentaci\'on r\'apida;
  \item integraci\'on para contratos y seguridad;
  \item E2E para validar experiencia real por rol;
  \item mutaci\'on para medir eficacia real de la suite.
\end{itemize}

\textbf{Resultado esperado}:
combinar profundidad y coste de mantenimiento,
evitar dependencia exclusiva de una sola capa de pruebas.

\subsection{Decisi\'on D14: incluir mutaci\'on (PIT/Stryker)}

\textbf{Raz\'on de fondo}:
la cobertura porcentual no garantiza detecci\'on efectiva de fallos.

\textbf{Valor aportado}:
los mutantes supervivientes ayudan a localizar huecos de aserciones
y mejorar calidad de tests en c\'odigo cr\'itico.

\section{Decisiones de operaci\'on y despliegue}

\subsection{Decisi\'on D15: GitHub Actions como columna vertebral CI/CD}

\textbf{Razones}:
\begin{itemize}
  \item integraci\'on nativa con repositorio;
  \item trazabilidad de checks por commit/PR;
  \item menor sobrecarga de operaci\'on para equipo reducido.
\end{itemize}

\subsection{Decisi\'on D16: despliegue en PaaS gestionado}

\textbf{Razones}:
\begin{itemize}
  \item minimiza administraci\'on de infraestructura;
  \item acelera iteraci\'on y validaci\'on de cambios;
  \item mantiene coste compatible con contexto acad\'emico.
\end{itemize}

En la implementaci\'on concreta se utiliza Render como PaaS gestionado para
el despliegue de los servicios y PostgreSQL gestionado. Para el canal de
correo se adopta Brevo mediante API HTTP, evitando depender de SMTP en
producci\'on, ya que algunos PaaS restringen puertos salientes asociados al
protocolo de correo. Inicialmente se utiliz\'o SendGrid como proveedor
de correo transaccional; tras la finalizaci\'on de su plan gratuito
se migr\'o a Brevo, que permite enviar correo a cualquier destinatario
en su plan gratuito.
El dise\~no adaptativo del \texttt{EmailService} soporta m\'ultiples
proveedores (Brevo, Resend, SendGrid) mediante un patr\'on desacoplado
que permite migrar de proveedor sin modificar l\'ogica de negocio.

\textbf{Trade-off}:
menor control fino de infraestructura frente a VPS/Kubernetes,
considerado aceptable para objetivos y alcance del TFG.

\section{Matriz consolidada de decisiones}

\begin{table}[ht]
\centering
\begin{tabular}{|P{1.6cm}|P{3.8cm}|P{3.2cm}|P{4.4cm}|}
\hline
\textbf{ID} & \textbf{Decisi\'on} & \textbf{Alternativa principal descartada} & \textbf{Raz\'on dominante} \\
\hline
D1 & Alcance MVP priorizado & Backlog completo desde inicio & Reducir riesgo de no entrega funcional \\
\hline
D3 & SPA + API desacoplada & Monolito SSR & Mantenibilidad y evoluci\'on de UX \\
\hline
D6 & Spring Boot + Java & Node.js/Django & Menor incertidumbre en seguridad/persistencia \\
\hline
D7 & React + TS + Vite & Angular/Vue & Equilibrio productividad-control tipado \\
\hline
D8 & JWT + refresh & Sesi\'on server-side & Ajuste natural a cliente API-first \\
\hline
D11 & Flyway + validate & ddl-auto update & Trazabilidad y control de esquema \\
\hline
D12 & Cloud en producci\'on & Solo almacenamiento local & Evitar p\'erdida en entornos ef\'imeros \\
\hline
D13 & Testing en capas + mutaci\'on & Solo unitarias/cobertura & Robustez real frente a regresiones \\
\hline
D16 & PaaS gestionado & VPS/Kubernetes & Tiempo y coste operativos contenidos \\
\hline
\end{tabular}
\caption{Resumen de decisiones clave y justificaci\'on dominante}
\label{tab:resumen-decisiones}
\end{table}

\section{An\'alisis coste-beneficio por bloques de decisiones}

Para reforzar la justificaci\'on,
se analiza cada bloque de decisiones en t\'erminos de coste inicial,
beneficio obtenido y retorno en el contexto del TFG.

\subsection{Bloque de arquitectura}

\textbf{Coste inicial}:
incremento de complejidad de integraci\'on entre frontend y backend,
definici\'on de contratos API y mayor superficie de pruebas.

\textbf{Beneficio}:
separaci\'on de responsabilidades,
facilidad para aislar incidencias,
evoluci\'on independiente de interfaces sin rehacer l\'ogica de negocio.

\textbf{Retorno}:
alto. El coste adicional de arranque se recupera
al reducir deuda t\'ecnica estructural y facilitar extensiones del sistema.

\subsection{Bloque de stack tecnol\'ogico}

\textbf{Coste inicial}:
tiempo de configuraci\'on de toolchain,
definici\'on de convenciones de proyecto y aprendizaje cruzado entre capas.

\textbf{Beneficio}:
ecosistema maduro,
documentaci\'on abundante,
integraci\'on directa con seguridad, testing y despliegue automatizado.

\textbf{Retorno}:
alto. La reducci\'on de incertidumbre durante la implementaci\'on
ha compensado sobradamente la configuraci\'on inicial.

\subsection{Bloque de seguridad}

\textbf{Coste inicial}:
dise\~no de flujo de tokens,
gesti\'on de expiraciones y tratamiento de errores de autenticaci\'on.

\textbf{Beneficio}:
control de acceso robusto,
alineaci\'on con arquitectura stateless,
capacidad de endurecimiento incremental.

\textbf{Retorno}:
muy alto. En un sistema acad\'emico con datos sensibles,
la inversi\'on en seguridad es estructural, no opcional.

\subsection{Bloque de calidad y pruebas}

\textbf{Coste inicial}:
esfuerzo significativo en dise\~no de suites,
mantenimiento de mocks y pipelines de verificaci\'on.

\textbf{Beneficio}:
disminuci\'on de regresiones,
detenci\'on temprana de defectos,
confianza objetiva para despliegue continuo.

\textbf{Retorno}:
muy alto. La calidad de salida mejora de forma medible
y se reduce el tiempo de correcci\'on en fases tard\'ias.

\subsection{Bloque de despliegue y operaci\'on}

\textbf{Coste inicial}:
configuraci\'on de entornos,
variables seguras y automatismos de publicaci\'on.

\textbf{Beneficio}:
repetibilidad,
trazabilidad de releases,
validaci\'on continua sobre entorno similar a producci\'on.

\textbf{Retorno}:
alto. La operativa se vuelve estable y menos dependiente
de ejecuciones manuales fr\'agiles.

\section{Riesgos asociados a decisiones y medidas de contingencia}

Ninguna decisi\'on de ingenier\'ia es neutra;
por ello se explicitan riesgos y mitigaciones por cada eje principal.

\begin{table}[ht]
\centering
\begin{tabular}{|P{3.4cm}|P{4.3cm}|P{4.5cm}|}
\hline
\textbf{Eje de decisi\'on} & \textbf{Riesgo principal} & \textbf{Mitigaci\'on aplicada} \\
\hline
Arquitectura desacoplada & Errores de contrato entre cliente y API & DTOs expl\'icitos, pruebas de integraci\'on y E2E \\
\hline
JWT + refresh & Gesti\'on incorrecta de expiraciones/revocaci\'on & Pol\'iticas de expiraci\'on, filtros dedicados y validaci\'on de contexto \\
\hline
Flyway como autoridad & Fallos de migraci\'on en despliegue & Versionado incremental, validaci\'on previa y backups de datos \\
\hline
Cloud en producci\'on & Dependencia de servicio externo & Configuraci\'on por entorno y estrategia fallback local \\
\hline
CI/CD automatizada & Falsos positivos por inestabilidad de tests & Separaci\'on por etapas, reintentos controlados y revisi\'on de flaky tests \\
\hline
\end{tabular}
\caption{Riesgos derivados de decisiones y medidas de contingencia}
\label{tab:riesgos-decisiones}
\end{table}

\section{Trazabilidad de decisiones con requisitos y evidencia}

Para que la justificaci\'on no quede en el plano discursivo,
se vinculan decisiones con requisitos y con evidencia emp\'irica observada
en implementaci\'on y verificaci\'on.

\begin{table}[ht]
\centering
\begin{tabular}{|P{1.6cm}|P{2.8cm}|P{2.9cm}|P{4.7cm}|}
\hline
\textbf{ID} & \textbf{Requisito(s) relacionado(s)} & \textbf{Evidencia principal} & \textbf{Conclusi\'on de valor} \\
\hline
D3 & RQNF-004, RQNF-008 & Separaci\'on frontend/backend en repositorio y pipelines & Mejora mantenibilidad y portabilidad \\
\hline
D8 & RQNF-009, REGN-001 & Filtros JWT y validaciones de autorizaci\'on en servicios & Refuerza control de acceso por rol/contexto \\
\hline
D11 & RQNF-001, RQNF-004 & Hist\'orico de migraciones versionadas & Reduce riesgo de deriva de esquema \\
\hline
D13 & RQNF-001, RQNF-006 & Suites unitarias/integraci\'on/E2E en verde y mutaci\'on alta & Mayor fiabilidad y menor regresi\'on \\
\hline
D16 & RQNF-008, RQNF-004 & Workflows CI/CD y despliegue automatizado en PaaS & Operaci\'on reproducible con bajo coste \\
\hline
\end{tabular}
\caption{Trazabilidad entre decisiones, requisitos y evidencia}
\label{tab:trazabilidad-decisiones}
\end{table}

\section{Decisiones descartadas y por qu\'e no se adoptaron en esta fase}

\subsection{Autenticaci\'on federada completa en MVP}

No se adopt\'o como requisito de salida del MVP por impacto en integraci\'on,
auditor\'ia de seguridad y superficie de pruebas,
frente al valor incremental inmediato sobre el n\'ucleo funcional.

\subsection{2FA integral desde primera versi\'on}

Se consider\'o relevante,
pero su adopci\'on temprana pod\'ia desviar esfuerzo de procesos acad\'emicos
cr\'iticos (entregas/evaluaciones) dentro del horizonte temporal disponible.

\subsection{Orquestaci\'on avanzada de infraestructura}

Kubernetes y configuraciones equivalentes se descartaron
por sobrecoste operativo para un equipo reducido y un contexto de TFG,
siendo m\'as adecuado reservarlo para fases de escalado posterior.

\section{Conclusiones del cap\'itulo}

La justificaci\'on detallada muestra que las decisiones del proyecto
no se basaron en preferencia tecnol\'ogica aislada,
sino en una evaluaci\'on sistem\'atica de valor, riesgo, coste y viabilidad.

El resultado es una arquitectura coherente con el ERS/DAS,
con compromisos expl\'icitos y argumentos t\'ecnicos defendibles ante tribunal,
que proyecta una hoja de ruta de evoluci\'on realista para fases futuras.
